\section{Expectation}

The expectation is the integral of a random variable over the probability space with respect to the probability measure.

Two things stand out about independent random variables.

\begin{enumerate}
    % [CORRECTED] was: "If $\eta_1, \eta_2, \eta_3, \ldots$ and $f, g, h, \ldots$ are measurable functions" --- every
    % random variable is measurable, so as written the claim is that any family at all has independent images;
    % $\eta_1 = \eta_2$ a single coin flip and $f = g = \id$ refutes it.
    \item If $\eta_1, \eta_2, \eta_3, \ldots$ are independent and $f, g, h, \ldots$ are measurable functions, then $f\of{\eta_1}, g\of{\eta_2}, h\of{\eta_3, \eta_4}, f\of{\eta_5}, \ldots$ are all independent.

    \item If $\eta_1, \ldots, \eta_n$ are independent, then
    \begin{align*}
        \E\of{\prod_{i=1}^{n} \eta_i} = \prod_{i=1}^{n} \E\of{\eta_i}
    \end{align*}
\end{enumerate}

\subsection{Conditional Expectation}

Often we learn something partway through and want the expectation of what is left. That is an expectation again, taken over a smaller sample space.

% [CORRECTED] was: "Given an event $A$ and a random variable $X$" --- at $\Pr(A) = 0$ the display below is $\frac{0}{0}$
% (every singleton in a null event is null, so the numerator vanishes too), so it defines nothing there, and the
% second display needs the same hypothesis, $\Pr\of{X = r \mid A}$ dividing by $\Pr(A)$ in its turn.
Given an event $A$ with $\Pr(A) \neq 0$ and a random variable $X$, we can define
\begin{align*}
    \E\of{X \mid A} = \frac{1}{\Pr(A)} \sum_{x \in A} X(x) \cdot p(x)
\end{align*}
where $p(x)$ is the density function. It is easy to show that
\begin{align*}
    \E\of{X \mid A} = \sum_{r \in X[\Omega]} r \cdot \Pr\of{X = r \mid A}
\end{align*}
% [CORRECTED] was: "$y_0 \mapsto \E\of{X \mid Y = y_0}$" in the display below --- that map has domain $\R$, not $\Omega$,
% so it is not a random variable in the sense of 1.1, and the outer $\E$ of the tower
% property has nothing to integrate against $\Pr$. Composing with $Y$ is what the next sentence claims has been done.
Given \textit{two} random variables $X$ and $Y$, we denote by $\E\of{X \mid Y}$ the \textbf{random variable}
\begin{align*}
    \omega \mapsto \E\of{X \mid Y = Y(\omega)}
\end{align*}

The point of packaging it as a random variable rather than as a list of numbers is that it can then be averaged again.

\begin{boxtheorem}[Tower Property of Conditional Expectation]
    For all random variables $X$ and $Y$,
    \begin{align*}
        \E\of{\E\of{X \mid Y}} = \E(X)
    \end{align*}
\end{boxtheorem}

The above type-checks because $\E\of{X \mid Y}$ is a random variable.

% Not from the lecture: the lecture defined the notation for one conditioning variable and then used it for several; the extension below was supplied.
Several random variables at once are no different: $\E\of{X \mid Y_1, \, \ldots, \, Y_k}$ is the random variable $\omega \mapsto \E\of{X \mid Y_1 = Y_1(\omega), \, \ldots, \, Y_k = Y_k(\omega)}$, conditioning on the one event that records what all of them did. Both this and $\E\of{X \mid Y}$ need each of those events to have nonzero probability, which on a finite $\Omega$ whose outcomes all have positive probability---the only case we use---every $\omega$ supplies.

\Cref{Ch4:Sec:Martingales} leans on all of this rather harder than anything before it does.
